\newcommand{\GaussVTheoryGathering}[1]{\gdef\GaussVTheoryGatheringText{#1}}
\newcommand{\GaussVTheoryEarlyFolio}[1]{%
  \GaussVTheoryGathering{}%
  \fancypagestyle{GaussVTheory#1}{%
    \fancyhead{}%
    \ifodd#1\relax
      \fancyhead[C]{\small DES ERDMAGNETISMUS\ifnum#1=135\relax\else\ifnum#1=165\relax\else.\fi\fi}%
      \fancyhead[R]{\small #1}%
    \else
      \fancyhead[L]{\small #1}%
      \fancyhead[C]{\small ALLGEMEINE THEORIE}%
    \fi
    \ifnum#1=123\relax
      \fancyfoot[R]{\small 16\textsuperscript{*}\hspace{.18\headwidth}}%
    \else
      \fancyfoot[R]{\small \GaussVTheoryGatheringText\hspace{.10\headwidth}}%
    \fi
    \renewcommand{\headrulewidth}{0pt}%
  }%
  \thispagestyle{GaussVTheory#1}%
}
\fancypagestyle{GaussVTheoryOpening}{%
  \fancyhead{}%
  \fancyfoot[R]{\small 16\hspace{.18\headwidth}}%
  \renewcommand{\headrulewidth}{0pt}%
}
\SourceLeaf{125}\thispagestyle{GaussVTheoryOpening}
\fancyhead[L]{\small Allgemeine Theorie des Erdmagnetismus}
\renewcommand{\Article}[1]{\par\medskip\Needspace{4\baselineskip}\pdfbookmark[1]{Art. #1}{erdmagnetismus-art-#1}\centerline{\textbf{#1.}}\nobreak\smallskip}
\pdfbookmark[0]{Allgemeine Theorie des Erdmagnetismus}{erdmagnetismus-body}
\vspace*{1.38in}
\begin{center}
ALLGEMEINE THEORIE\par\medskip
{\small DES\par}\medskip
{\Large\addfontfeatures{LetterSpace=16}ERDMAGNETISMUS.\par}
\medskip
\rule{.09\linewidth}{.35pt}\par
\end{center}
Der rastlose Eifer, womit man in neuerer Zeit in allen Theilen der Erdoberfläche die Richtung und Stärke der magnetischen Kraft der Erde zu erforschen strebt, ist eine um so erfreulichere Erscheinung, je sichtbarer dabei das rein wissenschaftliche Interesse hervortritt. Denn in der That, wie wichtig auch für die Schifffahrt die möglichst vollständige Kenntniss der Abweichungslinie ist, so erstreckt sich doch ihr Bedürfniss eben nicht weiter, und was darüber hinausliegt, bleibt für jene beinahe gleichgültig. Aber die Wissenschaft, wenn gleich gern auch dem materiellen Interesse förderlich, lässt sich nicht auf dieses beschränken, sondern fordert für Alle Elemente ihrer Forschung gleiche Anstrengung.

Die Ausbeute der magnetischen Beobachtungen pflegt man auf den Erdkarten durch drei Systeme von Linien darzustellen, die man wohl die isogonischen, isoklinischen und isodynamischen Linien genannt hat. Diese Linien ändern ihre Gestalt und Lage im Laufe der Zeit sehr bedeutend, so dass Eine Zeichnung nur den Zustand der Erscheinung für einen bestimmten Zeitpunkt angibt. \textsc{Halley}'s Declinationskarte ist sehr verschieden von \textsc{Barlow}'s Darstellung im Jahr 1833; und \textsc{Hansteen}'s Inclinationskarte für 1780 weicht schon sehr stark von der jetzigen Lage der isoklinischen Linien ab: die Versuche, die Intensität darzustellen, sind noch zu neu, als dass sich bei derselben schon jetzt ähnliche Änderungen \par
\clearpage
\SourceLeaf{126}\GaussVTheoryEarlyFolio{122}
\noindent nachweisen liessen, die ohne Zweifel im Laufe der Zeit nicht ausbleiben werden. Alle diese Karten sind jetzt noch mehr oder weniger lückenhaft, oder theilweise unzuverlässig: es steht aber zu hoffen, dass, wenn sie auch die Vollständigkeit, wegen der Unzugänglichkeit einiger Theile der Erdfläche, nicht ganz erreichen können, sie doch mit raschen Schritten sich ihr mehr nähern werden.

Vom höhern Standpunkt der Wissenschaft aus betrachtet ist aber diese möglichst vollständige Zusammenstellung der Erscheinungen auf dem Wege der Beobachtung noch nicht das eigentliche Ziel selbst: man hat damit nur ähnliches gethan, wie der Astronom, wenn er z. B. die scheinbare Bahn eines Kometen auf der Himmelskugel beobachtet hat. Man hat nur Bausteine, kein Gebäude, so lange man nicht die verwickelten Erscheinungen Einem Princip unterwürfig gemacht hat. Und wie der Astronom, nachdem das Gestirn sich seinen Augen entzogen hat, sein Hauptgeschäft erst anfängt, gestützt auf das Gravitationsgesetz aus den Beobachtungen die Elemente der wahren Bahn berechnet, und dadurch sogar sich in den Stand setzt, den weitern Lauf mit Sicherheit anzugeben: so soll auch der Physiker sich die Aufgabe stellen, wenigstens in so weit die ungleichartigen und zum Theil weniger günstigen Umstände es verstatten, die die Erscheinungen des Erdmagnetismus hervorbringenden Grundkräfte nach ihrer Wirkungsart und nach ihren Grössenwerthen zu erforschen, die Beobachtungen, so weit sie reichen, diesen Elementen zu unterwerfen, und dadurch selbst wenigstens mit einem gewissen Grade von sicherer Annäherung die Erscheinungen für die Gegenden, wohin die Beobachtung nicht hat dringen können, zu anticipiren. Es ist jedenfalls gut, dies höchste Ziel vor Augen zu haben, und die Gangbarmachung der dazu führenden Wege zu versuchen, wenn auch gegenwärtig, bei der grossen Unvollkommenheit des Gegebenen, mehr als eine entfernte Annäherung zu dem Ziele selbst noch nicht möglich ist.

Es ist nicht meine Absicht, hier diejenigen frühern erfolglosen Versuche zu erwähnen, wobei man ohne alle physikalische Grundlage das grosse Räthsel errathen zu können meinte. Eine physikalische Grundlage kann man nur solchen Versuchen zugestehen, welche die Erde wie einen wirklichen Magnet betrachten, und die erwiesene Wirkungsart eines Magneten in die Ferne allein der Rechnung unterstellen. Aber alle bisherigen Versuche dieser Art haben das gemein, dass man, anstatt zuerst zu untersuchen, \textit{wie} dieser grosse Magnet beschaffen sein müsse, um den Erscheinungen Genüge zu leisten, gleich gefasst darauf, eine \par
\clearpage
\SourceLeaf{127}\GaussVTheoryEarlyFolio{123}
\noindent einfache oder eine sehr zusammengesetzte Beschaffenheit hervorgehen zu sehen, vielmehr von vorne her von einer bestimmten einfachen Beschaffenheit ausging, und probirte, ob die Erscheinungen sich mit solcher Hypothese vertrügen. Indessen wiederholt sich hierin nur, was die Geschichte der Astronomie und der Naturwissenschaften von den Anfängen so vieler unserer Kenntnisse berichtet.

Die einfachste Hypothese dieser Art ist die, nur einen einzigen sehr kleinen Magnet im Mittelpunkt der Erde anzunehmen, oder vielmehr (da schwerlich jemand im Ernste an das wirkliche Vorhandensein eines solchen Magnets geglaubt hat) vorauszusetzen, der Magnetismus sei in der Erde so vertheilt, dass die Gesammtwirkung nach aussen der Wirkung eines fingirten unendlich kleinen Magnets äquivalire, ungefähr eben so, wie die Gravitation gegen eine homogene Kugel der Anziehung einer gleich grossen, im Mittelpunkt concentrirten Masse gleichkommt. In dieser Voraussetzung sind die beiden Punkte, wo die Fortsetzung der magnetischen Axe jenes Centralmagnets die Erdfläche schneidet, die magnetischen Pole der Erde, in denen die Magnetnadel vertical steht, und zugleich die Intensität am grössten ist; in dem grössten Kreise mitten zwischen beiden Polen (dem magnetischen Aequator) wird die Inclination $=0$ und die Intensität halb so gross als in den Polen; zwischen dem magnetischen Aequator und einem Pole hängt sowohl Inclination als Intensität nur von dem Abstande von jenem Aequator (der magnetischen Breite) ab, und zwar so, dass die Tangente der Inclination der doppelten Tangente dieser Breite gleich ist; endlich fällt die Richtung der horizontalen Nadel überall mit der Richtung eines nach dem nordlichen magnetischen Pole gezogenen grössten Kreises zusammen. Mit allen diesen nothwendigen Folgen jener Hypothese stimmt aber die Natur nur in roher Annäherung überein; in der Wirklichkeit ist die Linie verschwindender Inclination kein grösster Kreis, sondern eine Linie von doppelter Krümmung; bei gleichen Neigungen findet man nicht gleiche Intensitäten; die Richtungen der horizontalen Nadel sind weit davon entfernt, alle nach Einem Punkte zu convergiren u. s. f. Es reicht also schon die oberflächliche Betrachtung hin, die Verwerflichkeit dieser Hypothese zu zeigen: gleichwohl wendet man den Einen der obigen Sätze noch jetzt als eine Näherung an, um die Lage der Linie verschwindender Inclinationen aus solchen Beobachtungen abzuleiten, die in einiger Entfernung von ihr, bei mässigen Inclinationen, gemacht sind.

Von einer ähnlichen Hypothese war bereits vor 80 Jahren \textsc{Tobias Mayer} \par
\clearpage
\SourceLeaf{128}\GaussVTheoryEarlyFolio{124}
\noindent ausgegangen, nur mit der Modification, dass er den unendlich kleinen Magnet nicht in den Mittelpunkt der Erde, sondern etwa um den siebenten Theil des Erdhalbmessers davon entfernt setzte: doch behielt er, vermuthlich um grössere Verwicklung der Rechnung zu vermeiden, die an sich ganz willkürliche Beschränkung bei, dass die gegen die Axe des Magnets senkrechte Ebene durch den Mittelpunkt der Erde gehe. Auf diese Art fand er, bei einer freilich nur sehr kleinen Anzahl von Oertern, die beobachteten Abweichungen und Neigungen mit seiner Rechnung ganz gut übereinstimmend. Eine ausgedehntere Prüfung würde aber bald gezeigt haben, dass man mit jener Hypothese das Ganze der Erscheinungen dieser beiden Elemente nicht viel besser darstellen kann, als mit der zuerst erwähnten. Intensitätsbestimmungen gab es bekanntlich damals noch gar nicht.

\expandafter\def\csname HistoricalNoteText14\endcsname{Bei der Declination kommt sogar einmal ein Unterschied von mehr als 29 Grad vor: allein es ist billig, den Fehler der Rechnung nicht nach der Zahl der Declinationsgrade, sondern nach der wirklichen Ungleichheit zwischen der berechneten und beobachteten ganzen Richtung zu schätzen, wo er bei dem in Rede stehenden Orte $11\frac{1}{2}$ Grad beträgt.}
\textsc{Hansteen} ist einen Schritt weiter gegangen, indem er die Hypothese \textit{zweier} unendlich kleiner Magnete von ungleicher Lage und Stärke den Erscheinungen anzupassen versucht hat. Die entscheidende Prüfung der Zulässigkeit oder Unzulässigkeit einer Hypothese bleibt immer die Vergleichung der in ihr erhaltenen Resultate mit den Erfahrungen. \textsc{Hansteen} hat die seinige mit den Beobachtungen an 48 verschiedenen Oertern verglichen, unter denen sich jedoch nur 12 befinden, wo die Intensität mit bestimmt ist, und überhaupt nur 6, wo alle drei Elemente vorkommen. Wir treffen hier noch Differenzen zwischen der Rechnung und Beobachtung an, die bei der Inclination fast auf 13 Grad steigen\HistoricalNote{14}.

Wenn man nun so grosse Abweichungen den Forderungen nicht ent\-sprechend findet, die an eine genügende Theorie gemacht werden müssen, so kann man nicht umhin, den Schluss zu ziehen, dass die magnetische Beschaffenheit des Erdkörpers keine solche ist, für welche eine Concentrirung in Einen oder ein Paar einzelne unendlich kleine Magnete als Stellvertreterin gelten könnte. Es wird damit nicht geleugnet, dass mit einer \textit{grössern} Anzahl solcher fingirten Magnete zuletzt eine genügende Uebereinstimmung erreichbar werden könnte: allein eine ganz andere Frage ist, ob eine solche Form der Auflösung der Aufgabe gerathen sein würde; es scheint in der That, dass die schon bei zwei Magneten so
